\section{Overview}
\label{sec:overview}

This section offers a brief primer on Royco's structuring substrates: Dawn and Dusk.
Although this document specifies Dusk, the two substrates share their
core machinery: tranches, the per-tranche NAVs, and the profit and loss
waterfall.

\subsection{Tranches and NAVs}
\label{sec:overview:tranches}

A Royco market structures an asset into two complementary investment opportunities
with distinct risk/return profiles called \emph{tranches}:

\begin{itemize}
  \item \textbf{Senior tranche (ST):} Protected capital. Lower returns for receiving first-loss coverage.
  \item \textbf{Junior tranche (JT):} Levered capital. Higher returns for providing first-loss coverage.
\end{itemize}

Moreover, each tranche holds a perpetual obligation to the other. The
senior tranche pays the junior a share of its yield (risk premium)
for absorbing losses first. The junior tranche, in turn, stands
as a first-loss buffer: senior losses consume junior capital in full
before the senior tranche incurs any. Accounting for these obligations
requires distinguishing what each tranche \emph{deposited} from what it is
\emph{entitled to}. Hence, Royco tracks two NAVs (net asset value) per
tranche:

\begin{itemize}
  \item \textbf{Raw NAV:} The mark-to-market value of the tranche's \emph{deposited} assets.
  \item \textbf{Effective NAV:} The value that each tranche \emph{is entitled to} after its obligations to the other are settled.
\end{itemize}

\subsection{The profit and loss waterfall}
\label{sec:overview:waterfall}

Every accounting synchronization marks the tranche NAVs to market in three steps:

\begin{enumerate}
  \item \textbf{Measure:} Read the current raw NAVs and
    compute the delta against their last checkpointed values. Each delta
    represents the unrealized profit or loss of that tranche's holdings
    since the previous accounting sync.
  \item \textbf{Attribute:} Past coverage and risk premiums can leave a
    tranche entitled to more than its own holdings are worth. The excess
    ($\STeff - \STraw$ when positive, and symmetrically for JT) is a
    claim on the \emph{other} tranche's holdings. By conservation the two
    excesses are equal and opposite, so one tranche's excess is exactly
    the other's shortfall and the cross-claim runs one way. Each side's
    raw NAV delta is split pro rata over the claims on it. The result is
    a pair of effective NAV deltas that sum to the total profit and loss.
  \item \textbf{Settle:} The effective deltas are routed through the
    \emph{PnL waterfall} in fixed priority order:
    \begin{itemize}
      \item \emph{Loss:} Utilize JT's coverage buffer ($\JTeff$) before ST books
        any. Any coverage that JT advances to ST is booked as \emph{JT
        impermanent loss} (ST's debt to JT). Losses exceeding JT's
        buffer reduce $\STeff$ and are booked as \emph{ST impermanent
        loss} (JT's debt to ST). JT IL is transient: it exists only
        during the market's fixed-term regime, a window that lets the
        tranched asset recover before the junior realizes losses.
      \item \emph{Recovery:} Repay any debts before gains accrue.
        ST's IL has the first claim on appreciation from either side,
        and JT's IL has the second claim on senior appreciation only.
      \item \emph{Yield:} Junior appreciation remaining after ST IL
        repayment accrues to JT in full. Senior appreciation remaining
        after all IL repayment is true yield, split between ST and JT
        at the rate set by the yield distribution model (YDM).
    \end{itemize}
\end{enumerate}

We denote the attribution and settlement steps as $F$ throughout this
document. Holding the checkpoint fixed, $F$ maps the current raw NAVs
to the new effective NAVs:
\[
  F\colon (\STraw,\ \JTraw;\ \text{checkpoint}) \longmapsto (\STeff,\ \JTeff).
\]
The semicolon separates variables from parameters: fixing the
checkpoint yields a function of the raw NAVs alone. The checkpoint holds the previous sync's raw NAVs,
effective NAVs, outstanding ILs, market state, and yield-accrual
bookkeeping (the time-weighted YDM accumulator and its timestamps). It
is read from storage when the sync begins and cannot change while it
runs. Every property stated for $F$ holds for each fixed checkpoint.

$F$ holds two key properties:

\begin{enumerate}
  \item \textbf{NAV conservation:} the inputs sum to the outputs,
    $\STraw + \JTraw = \STeff + \JTeff$, for every checkpoint that
    itself conserves. Every stored checkpoint conserves, because it is
    itself either the conserving genesis checkpoint seeded at market
    creation or a prior output of $F$. $F$ redistributes value between the tranches without
    creating or destroying it.
  \item \textbf{Piecewise-affine:} $F$ is built from chains of $\min$,
    $\max$, and affine pieces with slopes $\in [0, 1]$, so it is
    continuous and non-decreasing in each input. NAV conservation then
    makes it 1-Lipschitz in each input. The existence proof invokes the
    junior direction (\cref{sec:existence}).
\end{enumerate}

\subsection{Dawn to Dusk}
\label{sec:overview:dawn-to-dusk}

Dawn transforms the \emph{risk profile} of any asset, but not its
\emph{liquidity profile}: senior shares have no exit path other than
primary redemption.

Dusk extends the junior tranche to transform both. In Dawn, junior
capital may sit in any priceable asset (Treasuries, a money market
fund, or co-investment alongside the senior tranche), provided its
value is exogenous to the market itself. Dusk drops that restriction:
junior capital is deployed into an AMM pool paired against the senior
tranche's shares.

\begin{itemize}
  \item \textbf{Risk transformation:} Coverage still flows from JT to ST exactly
    as in Dawn: the pool position is JT's first-loss collateral.
  \item \textbf{Liquidity transformation:} The pool position gives senior
    holders an instant secondary market, constantly quoting a price to
    buy and sell tranche shares for stables.
\end{itemize}

For the purposes of this document, the Dusk JT's asset is a Balancer V3
Pool Token (BPT) whose constituents are the senior tranche share and a
stable quote asset (USDC, sUSDS, or srRoyUSDC).

\textbf{The recursion:} Marking JT to market is no longer a plain oracle
read. To value the BPT, the pool's senior share leg must be priced at
the senior rate $y = \STeff / N$, where $N$ is the total supply of senior shares.
But $\STeff$ is not an observable quantity: it is the output of the
sync map $F$, and $F$ cannot run without its $\JTraw$ input, which is
the BPT value we set out to compute. The dependency chain closes on
itself, as \cref{fig:recursion} shows.

\begin{figure}[!ht]
\centering
\begin{tikzpicture}[
  ->, >={Stealth[length=2.2mm]},
  font=\small,
  shorten >=4pt, shorten <=4pt,
]
  \node (y)  at (90:1.6)  {$y$};
  \node (jt) at (-30:1.6) {$\JTraw$};
  \node (st) at (210:1.6) {$\STeff$};
  \draw (y)  to[bend left=30] node[right=8pt]  {Value the BPT} (jt);
  \draw (jt) to[bend left=30] node[below=8pt]  {Apply $F$} (st);
  \draw (st) to[bend left=30] node[left=8pt]   {Value the share} (y);
\end{tikzpicture}
\caption{The recursion. Valuing the junior collateral needs the senior
rate, and the senior rate depends on that valuation.}
\label{fig:recursion}
\end{figure}

Economically, the junior tranche's collateral contains the very senior
shares its coverage protects. No evaluation order can resolve the cycle:
the two unknowns, $\JTraw$ and $\STeff$ ($y$ is simply $\STeff$ per
share), must be determined simultaneously. Under conditions made
precise in \cref{sec:existence}, the values of $\STeff$ that satisfy
both dependencies at once form a contiguous block: price the BPT at its
implied rate $y = \STeff/N$, apply $F$, and the same $\STeff$ returns.
Each such self-reproducing value is a \emph{fixed point}, and the
published value is the largest. \Cref{sec:setup}
formulates the equation, \cref{sec:existence} proves the published
value exists and is uniquely selected, and \cref{sec:solver} finds it
on-chain.

\subsection{Components of a Dusk market}
\label{sec:overview:pieces}

A Dusk market adds three components to the Dawn stack (pool, kernel,
hook) and reuses the Dawn accountant. Each plays one role in resolving
the recursion:

\begin{itemize}
  \item \textbf{Gyro E-CLP pool:} A Balancer V3 elliptic concentrated
    liquidity pool over the senior share and quote asset. The senior
    leg is priced by the rate the kernel publishes, so the solved rate
    is what the pool uses to value senior shares.
  \item \textbf{Dusk kernel:} One contract, two paths. Write: the sync
    entrypoint solves the fixed point (\cref{sec:solver}), commits the
    accountant checkpoint, and caches the solved rate $y^\star$. Read:
    the kernel reports the cached rate $y^\star$ to the pool instantly.
  \item \textbf{Balancer V3 Dusk hook:} Forwards Balancer's before-operation
    callbacks (swap, add and remove liquidity) into the kernel's sync
    entrypoint. Balancer reloads balances and rates after the hook so
    every standard operation executes at a fresh senior share rate.
  \item \textbf{Accountant:} Shared with Dawn. During a solve, the
    kernel calls its view-only preview repeatedly to evaluate $F$ at
    candidate rates. After convergence, exactly one call commits the
    resultant checkpoint.
\end{itemize}
